\section{Introduction}

Industrial fault diagnosis \cite{precup2015overview}, which identifies the underlying root causes of observed deviations in manufacturing and process equipment, is a critical technology in the era of Industry 4.0 \cite{precup2015overview}. Over the past decade, data-driven approaches have gained considerable prominence because they offer rapid and accurate inferences without relying on complex, first-principles physical models. Traditionally, data-driven fault diagnosis \cite{zhao2021semisupervised} has been formulated as a supervised machine learning classification problem \cite{zhao2021semisupervised,du2023fault} \cite{wang2020deep}. In scenarios where abundant, well-labeled historical data is available, intelligent diagnostic models—such as support vector machines, random forests, and deep convolutional neural networks \cite{he2016deep}—can identify specific fault categories with exceptional accuracy \cite{wang2020deep}. However, modern industrial environments impose strict reliability and safety constraints. It is generally impermissible to intentionally operate safety-critical machinery \cite{lee2014prognostics}, such as steam turbines, high-pressure chemical reactors, or electrical transmission grids, in a malfunctioning state merely to collect fault data \cite{wang2020deep}. Furthermore, highly durable components \cite{zhang2022intelligent}, such as premium bearings and gearboxes, may operate for years before experiencing significant degradation, making comprehensive data collection prohibitively time-consuming \cite{wang2020deep}. Consequently, a scarcity of labeled fault data is a prevalent and substantial barrier in practice. This scarcity fundamentally limits standard supervised diagnostic models, which inevitably fail when confronted with novel or rare fault types that lack historical training instances.

To address the critical challenge of unavailable fault data, Zero-Shot Learning (ZSL) has been introduced into the field of fault diagnosis \cite{precup2015overview}, establishing the paradigm of Zero-Shot Fault Diagnosis (ZSFD) \cite{lampert2014attribute,xian2019zero}. The core principle of ZSFD is to robustly bridge the knowledge gap between "seen" fault classes—which have abundant training samples—and "unseen" fault classes—which have none. This bridge is typically constructed utilizing a shared semantic space designed by human domain experts. By establishing a mapping between the diagnostic signal feature space and the semantic space using data from seen classes, a zero-shot model can subsequently recognize unseen faults strictly by evaluating their semantic representations. Recently, generative-based zero-shot learning frameworks, notably Generative Adversarial Networks (GANs) \cite{goodfellow2014generative} and Variational Autoencoders (VAEs) \cite{lv2020hybrid}, have demonstrated significant potential in augmenting these datasets \cite{wang2020deep}. These data-driven generative models learn to synthesize discriminative pseudo-samples or feature embeddings for unseen fault classes conditioned on their specific semantic descriptions. For example, recent methodologies such as CycleGAN-SD \cite{liao2024cyclegan} transform fault features across domains based on calculated semantic distances, alleviating the extreme data imbalance and allowing standard classifiers to be trained directly on the synthesized unseen target samples \cite{wang2020deep}. 

Despite these promising computational advances, existing generative zero-shot fault diagnosis \cite{precup2015overview} methodologies suffer from several critical limitations. First, nearly all current models rely heavily on manual, discrete binary attributes (e.g., rigid $0/1$ matrices) to define the shared semantic space \cite{feng2021fault}. While binary attributes effectively capture basic, mutually exclusive states (e.g., "valve fully closed" versus "valve functioning normally"), they severely lack the representational capacity to convey rich, fine-grained semantic information, such as progressive degradation dynamics, subtle thermal context, or complex causal relations inherent in process faults. This hard-coded discretization strips away the continuous, nuanced nature of real-world industrial fault descriptions, fundamentally limiting the model's capacity to differentiate between highly similar unseen failure modes. Second, the predominant generative engines utilized in current ZSFD studies—specifically GANs—suffer from well-documented training instabilities and mode collapse vulnerabilities \cite{arjovsky2017wasserstein} \cite{wang2020deep}. The adversarial Min-Max optimization strategy is notoriously difficult to converge, and the generated signals often lack the essential diversity required to faithfully represent the true distribution of an unseen fault class. Furthermore, these architectures frequently experience a severe "domain shift," \cite{tzeng2017adversarial,long2015learning} wherein the generated features of unseen classes remain heavily biased toward the seen domains they were originally trained on, resulting in degraded fault classification accuracy during practical deployment \cite{wang2020deep}.

To mathematically formulate the challenge of cross-modal alignment and domain shift in existing approaches, let $x \in \mathbb{R}^D$ denote a physical diagnostic signal and $e_{\text{text}} \in \mathbb{R}^K$ represent its corresponding textual semantic embedding. Existing attribute regressor-based models typically enforce semantic consistency through a simple regression mapping, which struggles to align complex multimodal data distributions. The objective fundamentally requires a robust probabilistic transition that accurately models the conditional data distribution $p(x | e_{\text{text}})$, while minimizing the cross-modal discrepancy problem:
\begin{equation}
\label{eq:cross_modal_gap}
\mathcal{L}_{\text{gap}} = \mathbb{E}_{x \sim p_{\text{data}}, e_{\text{text}}} \left[ \| f_{\text{visual}}(x) - f_{\text{semantic}}(e_{\text{text}}) \|_2^2 \right]
\end{equation}
where $f_{\text{visual}}$ and $f_{\text{semantic}}$ denote modality-specific feature encoders. However, simply minimizing Euclidean distance in a latent space under unstable adversarial training fails to reliably generate high-fidelity, diverse industrial signals for entirely unseen classes.

Recognizing the stark limitations of discrete binary attributes and unstable adversarial networks, Denoising Diffusion Probabilistic Models (DDPMs) \cite{ho2020denoising,dhariwal2021diffusion} present a highly compelling alternative \cite{feng2021fault}. Diffusion probabilistic models offer superior training stability and can synthesize incredibly diverse and high-fidelity signal distributions by reversing a structured Markovian noise process. Concurrently, the advent of Large Language Models (LLMs) and Vision-Language architectures like CLIP (Contrastive Language-Image Pretraining) \cite{radford2021learning} has revolutionized cross-modal semantic understanding, enabling the extraction of continuous, high-dimensional textual embeddings that far exceed the expressive power of manually crafted binary matrices \cite{wang2020deep}.

To deeply address the shortcomings of current ZSFD approaches, this paper proposes a novel framework named the CLIP-Guided Dual-Path Diffusion Model (CDDM). Instead of relying on rigid binary parameters, CDDM utilizes LLM-based continuous semantic embeddings. By encoding comprehensive textual descriptions of fault causes and observable physical effects, the semantic space forms a continuous, dense representation that expertly bridges the cross-modal gap between textual semantics and high-frequency diagnostic sensing. To generate realistic diagnostic samples from these continuous embeddings, we abandon adversarial training altogether in favor of a highly stable generative diffusion process. Specifically, we design a dual-path conditional diffusion architecture composed of a 1D Convolutional Neural Network (CNN) feature path focused strictly on time-domain transient patterns, and a parallel Spectral/Attention path designated for capturing time-frequency harmonic behavior. This dual-path design explicitly models both instantaneous operational condition shifts and long-term cyclic degradations. To seamlessly integrate the continuous textual semantics into the reverse generation process, we employ a Cross-Attention conditioning mechanism. Finally, to strictly enforce semantic alignment and definitively eradicate the typical domain shift problem, a multimodal contrastive consistency loss is introduced directly into the diffusion training framework, effectively replacing traditional and often insufficient attribute regressors.

In summary, the main contributions of this article are highlighted as follows:
\begin{itemize}
    \item We propose a pioneering cross-modal zero-shot fault diagnosis \cite{precup2015overview} paradigm that successfully eliminates the traditional reliance on hard-coded discrete binary attributes. By leveraging LLM-driven continuous semantic embeddings, our approach rigorously preserves fine-grained, textual descriptions of industrial faults.
    \item We develop a generative Dual-Path Diffusion framework (incorporating both time and time-frequency domains) that synthesizes highly diverse and robust zero-shot fault signals, effectively overcoming the instability and mode collapse vulnerabilities inherent to GAN-based prior art.
    \item We introduce a novel multi-modal contrastive consistency loss. This mechanism forces the generated physical diagnostic signals to precisely align with their guiding continuous text semantics, effectively bridging the cross-modal gap and drastically mitigating the domain shift phenomenon in unseen class generation.
    \item We perform comprehensive testing across multiple challenging, industry-standard benchmarks, including the Three-Phase Transmission System (TPTS), the Tennessee Eastman Process (TEP), and a real-world Hydraulic System dataset. Extensive experimental data proves that the proposed CDDM achieves state-of-the-art zero-shot diagnostic performance compared against existing counterparts.
\end{itemize}